\documentclass[a4paper,twoside,12pt]{article}

\usepackage{amsmath}
\usepackage{amsfonts}
\usepackage{amssymb}
\usepackage{fancyhdr}
\usepackage{stmaryrd}
\usepackage{textcomp}
\usepackage[french]{babel}
\usepackage{variations}
\usepackage{graphicx}
\usepackage{array}
\usepackage{verbatim}
\usepackage{bbm}
\usepackage{ifthen}
\usepackage[paper=a4paper,hmargin=28mm,top=8mm,bottom=52.5mm,includehead=true,twoside=true,headsep=5mm,driver=xetex,twoside=true,bindingoffset=6mm]{geometry}

%textpos : permet de placer des éléments par leur coordonnées absolues sur la page
% utilisé pour les entêtes et pied de pages

\usepackage[absolute]{textpos}

\setlength{\TPHorizModule}{1mm}
\setlength{\TPVertModule}{1mm}

% appendix : permet d'écrire en couleur

\usepackage[usenames,dvipsnames,table]{xcolor}

%% pdfpages : utilisé pour insérer la page de garde avant le document

\usepackage{pdfpages}

% caption : centre les légendes des figures

\usepackage[center]{caption}

% wrapfig : permet de créer des figures entourées de texte

\usepackage{wrapfig}

%placeins : permet d'obliger Latex à insérer les figures en attente avant de continuer

\usepackage{placeins}

% subcaption : permet de créer des figures avec plusieurs sous-figures

\usepackage{subcaption}

% layouts : affichage d'une longueur

\usepackage{layouts}

% tabularx : mise en forme des tableaux
\usepackage{color}
\usepackage{longtable}
\usepackage{colortbl,arydshln}
\usepackage{tabularx}

% eurosym : symbole euro

\usepackage{eurosym}

\usepackage{tocvsec2}

% xelatex
\usepackage{xltxtra}

\usepackage{lastpage}

\usepackage{xparse}
\usepackage{hyperref}

\usepackage[nonumberlist]{glossaries}
\newcommand{\str}[1]{\renewcommand{\arraystretch}{#1}}

\makeatletter
\let\tobi@kept@section\section
\RenewDocumentCommand{\section}{s}
  {\IfBooleanTF{#1}{\tobi@kept@section*}{\tobi@section}}
\NewDocumentCommand{\tobi@section}{o m o}
  {\IfNoValueTF{#3}% No short title for the header
     {\IfNoValueTF{#1}% No title for the toc
        {\tobi@kept@section{#2}}
        {\tobi@kept@section[#1]{#2}}%
     }
     {\IfNoValueTF{#1}% No title for the toc
        {\tobi@kept@section[#2]{#2\sectionmark{#3}}\sectionmark{#3}}
        {\tobi@kept@section[#1]{#2\sectionmark{#3}}\sectionmark{#3}}%
    }%
  }
\makeatother


\newcommand{\refbiblio}[1]{\cite{#1}}
\newcommand{\norefbiblio}[1]{\nocite{#1}}

\newcommand{\mymark}[1]{\markboth{\MakeUppercase{#1}}{\MakeUppercase{#1}}}



\newenvironment{intro}
{
\section*{Introduction}
\addcontentsline{toc}{section}{Introduction}
\mymark{Introduction}
}
{

}

\newenvironment{remerciements}
{
\section*{Remerciements}
\addcontentsline{toc}{section}{Remerciements}
\mymark{Remerciements}
}
{

}

\newenvironment{conclusion}
{
\section*{Conclusion}
\addcontentsline{toc}{section}{Conclusion}
\mymark{Conclusion}
}
{

}

\newenvironment{resume}
{
\section*{Résumé}
\addcontentsline{toc}{section}{Résumé}
\mymark{Résumé}
}
{

}


\newenvironment{sous-partie}[1]
{
\subsection{#1}
}
{

}

\newenvironment{sous-sous-partie}[1]
{
\subsubsection{#1}
}
{

}

\newenvironment{liste}
{
\vspace{0.2cm}
\begin{list}{$\bullet$\hspace{0.3cm}}{\leftmargin=1.4cm}
}
{
\end{list}
\vspace{0.2cm}
}

\newenvironment{oliste}
{
\vspace{0.2cm}
\begin{enumerate}
}
{
\end{enumerate}
\vspace{0.2cm}
}



\newenvironment{partie}[1]
{
\section{#1}
}
{

}

\newcounter{numannexe}
\newenvironment{annexes}
{
\newpage
\thispagestyle{empty}
\mbox{}
\newpage
\addcontentsline{toc}{section}{Annexes}
\mymark{Annexes}
\vspace*{\fill}
\begin{center}
\Huge{\textbf{Annexes}}
\end{center}
\vspace*{\fill}
\newpage
\setcounter{numannexe}{0}

\renewenvironment{partie}[1]
{
	\refstepcounter{subsection}
	\subsection*{\arabic{subsection}. ##1}
}
{
}

}
{

}

\newenvironment{annexe}[1]
{
\setcounter{subsection}{0}
\refstepcounter{numannexe}
\addcontentsline{toc}{subsection}{Annexe \arabic{numannexe} : #1}
\mymark{Annexe \arabic{numannexe} : #1}
\section*{Annexe \arabic{numannexe} : #1}
}
{
	\FloatBarrier
}

\newenvironment{sous-annexe}[1]
{
\begin{subappendices}
\subsection{#1}
}
{
\end{subappendices}
}

% Permet de faire des cadres dans le document
\newcommand{\cadre}[1]{\renewcommand{\arraystretch}{0.4}\begin{array}{!{\vline}c!{\vline}}\hline\\#1\\\\\hline\end{array}}
\newcommand{\cadret}[1]{\renewcommand{\arraystretch}{0.4}\begin{tabular}{!{\vline}c!{\vline}}\hline\\#1\\\\\hline\end{tabular}}
\newcommand{\semicadre}[1]{\renewcommand{\arraystretch}{0.3}\begin{array}{c!{\vline}}#1\\\\\hline\end{array}}

% Commandes mathématiques

\newcommand{\equivalent}[2]{\!\!\renewcommand{\arraystretch}{1}\begin{array}[t]{c}\Huge\sim\\^{#1\rightarrow#2}\end{array}\!\!}
\newcommand{\tend}[2]{\!\renewcommand{\arraystretch}{1}\begin{array}[t]{c}-\!-\!\!\!\longrightarrow\\^{#1\rightarrow#2} \\ [-1.5ex]\end{array}\!}
\newcommand{\egal}[2]{\!\!\renewcommand{\arraystretch}{1}\begin{array}[t]{c}=\\^{#1\rightarrow#2}\end{array}\!\!}
\newcommand{\fin}{\vspace{0.2cm}\\}
\newcommand{\finq}{\vspace{0.5cm}\\}
\newcommand{\sh}{\mathrm{sh}\,}
\newcommand{\abs}[1]{\left\vert#1\right\vert }
\newcommand{\ch}{\mathrm{ch}\,}

\renewcommand{\o}[1]{\mathrm{o}\!\left(#1\right)}
\newcommand{\etoile}{\hspace*{1cm}$\star$\hspace*{0.5cm}}
\renewcommand{\th}{\mathrm{th}\,}
\renewcommand{\arcsin}{\mathrm{Arcsin}\,}
\renewcommand{\arccos}{\mathrm{Arccos}\,}
\newcommand{\argsh}{\mathrm{Argsh}\,}
\newcommand{\argch}{\mathrm{Argch}\,}
\newcommand{\argth}{\mathrm{Argth}\,}
\newcommand{\rg}{\mathrm{rg}\,}
\renewcommand{\arctan}{\mathrm{Arctan}\,}
\newcommand{\sq}{\hspace*{1.4cm}\stepcounter{sq}(\alph{sq})\hspace*{0.5cm}}
\newcommand{\rcl}{\begin{array}{rcl}}
\newcommand{\ea}{\end{array}}
\renewcommand{\tfrac}[2]{\textstyle\frac{#1}{#2}}
\newcommand{\Cl}[1]{$C^{#1}$}
\newcommand{\mathCl}[1]{C^{#1}}
\renewcommand{\t}[1]{\tilde{#1}}
\renewcommand{\l}{\lambda}
\newcommand{\ds}{\displaystyle}
\newcommand{\R}{\mathbb{R}}
\newcommand{\Q}{\mathbb{Q}}
\newcommand{\Z}{\mathbb{Z}}
\newcommand{\N}{\mathbb{N}}
\newcommand{\Ker}{\mathrm{Ker}\,}
\newcommand{\Vect}{\mathrm{Vect}}
\renewcommand{\lvert}{\left\vert}
\renewcommand{\Im}{\mathrm{Im}\,}
\renewcommand{\rvert}{\right\vert}
\newcommand{\mnk}{\mathcal{M}_n(K)}
\newcommand{\mnc}{\mathcal{M}_n(C)}
\newcommand{\ppcm}{\mathrm{ppcm}}
\newcommand{\Tr}{\mathrm{Tr}\,}
\renewcommand{\t}[1]{^t\!#1}
\newcommand{\scal}[2]{\left\langle #1|#2\right\rangle}
\newcommand{\scalindice}[4]{\phantom{\langle}_{#3}\!\left\langle #1|#2\right\rangle_{#4}}
\newcommand{\p}[1]{\left( #1 \right)}
\newcommand{\crochet}[1]{\left[ #1 \right]}
\newcommand{\nr}[1]{\left\|\,#1\,\right\|}
\newcommand{\tab}{\hspace*{1cm}}

\newcommand{\esp}[1]{\mathbb{E}\!\crochet{#1}}
\newcommand{\espcond}[2]{\mathbb{E}_{#1}\!\crochet{#2}}

\renewcommand{\P}[1]{\mathbb{P}\!\p{#1}}
\newcommand{\Pcond}[2]{\mathbb{P}_{#1}\!\p{#2}}

\renewcommand{\binom}[2]{\left(\begin{array}{c}#1\\#2\end{array}\right)}

\newcommand{\car}[1]{\mathbf{1}_{#1}}

\newcommand{\matdd}[4]{\left({\begin{array}{cc} #1 & #2\\ #3 & #4 \ea}\right)\vspace{0.05cm}}

\newcommand{\supp}[1]{\text{supp}\left(#1\right)}

\renewcommand{\ge}{\geqslant}
\renewcommand{\le}{\leqslant}

\newcommand{\lebesgue}{\mathcal{L}}

\newcommand{\Drond}{\mathcal{D}}

\renewcommand{\Re}{\mathcal{R}e}
\newcommand{\obs}[1]{\hat{#1}}
\newcommand{\ket}[1]{\vert #1 \rangle}
\newcommand{\bra}[1]{\langle #1 \vert}
\newcommand{\braindice}[2]{\!\phantom{\langle}_{#2}\!\left\langle #1\right\vert}

\newcommand{\prodscal}[2]{\left\langle#1,#2\right\rangle}

\newcommand{\vect}[2]{\left(\str{1}\begin{array}{cc}#1 & #2 \ea\right)}
\newcommand{\vecttrans}[2]{\left(\str{1}\begin{array}{c}#1 \\ #2 \ea\right)}

\renewcommand{\v}[1]{\underline{#1}}

\newcommand{\Dp}[2]{\dfrac{\partial #1}{\partial #2}}
\newcommand{\grad}{\mathrm{grad}\,}
\newcommand{\vgrad}{\v{\mathrm{grad}}\,}


\usepackage{listings}

\renewcommand{\thesection}{\Roman{section}}

\pagestyle{fancy}

\newlength{\headrulemargintop}
\setlength{\headrulemargintop}{10.15mm}


\newlength{\headbottommargin}

\setlength{\headheight}{2cm}
\setlength{\footskip}{2.65cm}

\renewcommand{\headrulewidth}{0pt}
\renewcommand{\footrulewidth}{0pt}

\setlength{\headsep}{6mm}
\addtolength{\headrulemargintop}{-\headsep}
\setlength{\headbottommargin}{\headsep}
\addtolength{\headheight}{\headsep}
\addtolength{\headsep}{-4mm}
\addtolength{\headbottommargin}{-\headsep}

\renewcommand{\headrule}{\vspace*{-\headrulemargintop}}
\renewcommand{\headrule}{\vspace*{-\headrulemargintop}}

\renewcommand{\sectionmark}[1]{\markright{\thesection\ \MakeUppercase{#1}}}
\renewcommand{\subsectionmark}[1]{\markright{\thesubsection\ \MakeUppercase{#1}}}

\fancyhead{}
\fancyfoot{}

\definecolor{pantone}{rgb}{0.,0.26274509803922,0.37647058823529}

\newlength{\logoskyline}
\setlength{\logoskyline}{3cm}
\newlength{\logoskylineleft}
\setlength{\logoskylineleft}{16.965250619335cm}
%\setlength{\logoskylineleft}{21cm}
%\addtolength{\logoskylineleft}{-\logoskyline}
%\addtolength{\logoskylineleft}{-6mm}

\fancyhead[LO]{%
\color{pantone}
\small
\str{0.9}
\begin{tabular}{@{}l@{}}
{\bf \fontspec{Gill Sans Light} INF551 - Raisonnement assisté par ordinateur}\\
{\bf \fontspec{Gill Sans Light} Rapport de projet - Compilateur vérifié en Coq}
\end{tabular}\\\vspace*{3.1mm}
\str{0.9}\begin{tabular}{@{}l@{}}
{\selectfont \fontspec{Gill Sans Light} Dorian Nogneng}\\{\selectfont \fontspec{Gill Sans Light}  Quentin Fiard}
\end{tabular}
\vspace*{\headbottommargin}
\begin{textblock}{5}(12,8)
\includegraphics[height=2.05cm]{./logo_x.eps}
\end{textblock}
\begin{textblock}{200}(34,19)
\hrule width 130.2486772486772mm
\end{textblock}}

\fancyhead[LE]{%
\color{pantone}
\small
\str{0.9}
\begin{tabular}{@{}l@{}}
{\bf \fontspec{Gill Sans Light} INF551 - Raisonnement assisté par ordinateur}\\
{\bf \fontspec{Gill Sans Light} Rapport de projet - Compilateur vérifié en Coq}
\end{tabular}\\\vspace*{3.1mm}
\str{0.9}\begin{tabular}{@{}l@{}}
{\selectfont \fontspec{Gill Sans Light} Dorian Nogneng}\\{\selectfont \fontspec{Gill Sans Light}  Quentin Fiard}
\end{tabular}
\vspace*{\headbottommargin}
\begin{textblock}{5}(6,8)
\includegraphics[height=2.05cm]{./logo_x.eps}
\end{textblock}
\begin{textblock}{200}(28,19)
\hrule width 130.2486772486772mm
\end{textblock}}

\fancyfoot[L]{\fancyplain{}\rightmark}

\fancyfoot[R]{\thepage\ / \pageref{LastPage}}

\newcommand{\definition}[2]{{\bf #1} : #2}

\makeglossaries

\renewcommand{\glsnamefont}[1]{\makefirstuc{#1}}

%\newglossaryentry{wattpic}
%{
%  name=watt-pic,
%  description={Puissance nominale produite par un système photovoltaïque sous %l'illumination maximale du site considéré}
%}

\begin{document}
%\includepdf{img/page0.pdf}

\setcounter{page}{1}

\clearpage

\tableofcontents

\clearpage

\begin{intro}

\begin{paragraph}{Description du projet\vspace{0.2cm}\\}
Ce projet consisté en la mise au point et la certification en Coq d'un compilateur pour un mini-langage fonctionnel. Nous allons dans ce rapport présenter le langage considéré, la machine abstraite de destination et le fonctionnement général du compilateur développé.
\end{paragraph}

\end{intro}

\begin{partie}{Langage étudié}

\begin{sous-partie}{Indices de De Bruijn et manipulation de $\lambda$-termes}

\begin{paragraph}{}
Le langage étudié est un modèle très simple des langages fonctionnels. Les termes de ce langage sont définis par la grammaire suivante :
\end{paragraph}
$$\begin{array}{rcll}
t&::=&x&\mbox{variable}\\
&|&\lambda\,x\cdot t&\mbox{fonction associant t à x}\\
&|&t\;u&\mbox{application de la fonction t à u}
\ea$$
\begin{paragraph}{}
Pour ôter l'indétermination lié à l'$\alpha$-équivalence de termes, l'énoncé propose d'utiliser les indices de De Bruijn : chaque variable liée est remplacée par un entier représentant le nombre de lieurs $\lambda$ la séparant de l'abstraction à laquelle elle est liée. Un exemple de terme de ce langage avec indices de Bruijn est par exemple $\lambda\cdot\lambda\cdot(0\;1)$, où le 1 est associé au $\lambda$ le plus externe.
\end{paragraph}

\end{sous-partie}

\begin{sous-partie}{Formalisation en Coq}

\begin{paragraph}{}
Les termes avec indices de De Bruijn sont représentés par le type \verb|DBT| en Coq. Ce type est défini par l'inductif suivant :
\end{paragraph}\vspace{-0.8cm}
\begin{center}
\begin{tabular}{p{5cm}}
\begin{verbatim}
Inductive DBT:Type :=
| Var : nat -> DBT
| Fun : DBT -> DBT
| Appl : DBT -> DBT -> DBT.
\end{verbatim}
\end{tabular}
\end{center}\vspace{-1cm}
\begin{paragraph}{}
Un terme de type \verb|Var x| représente une variable d'indice de De Bruijn $x\in\N$, \verb|Fun t| est la fonction qui a un terme $u$ associe le terme $t[0\leftarrow u]$, et \verb|Appl t1 t2| représente le terme $t1\;t2$, application de $t1$ à $t2$.
\end{paragraph}

\end{sous-partie}

\begin{sous-partie}{Formalisation des termes clos sous forme d'un prédicat}

\begin{paragraph}{}
Comme suggéré dans l'énoncé, nous avons procédé en deux étapes pour définir le prédicat ``Terme clos" à partir de la définition précédente des termes avec indices de De Bruijn.
\end{paragraph}
\begin{paragraph}{}
Nous avons commencé par formaliser le prédicat $C[n](t)$ suivant en Coq : pour tout $t$ terme avec indices de De Bruijn, $t$ satisfait $C[n](t)$ si et seulement si toutes les variables libres de $t$ sont d'indices $<n$. Ce prédicat est dénommé \verb|n_free t n| dans le code Coq (\verb|n_free t n = true| si et seulement si $C[n](t)$ est satisfait).
\end{paragraph}
\begin{paragraph}{}
Pour calculer ce prédicat, on commence par définir un prédicat auxiliaire\\\verb|n_k_free : DBT -> nat -> nat|, tel que \verb|n_k_free t n k = true| si et seulement si
toutes les variables libres de $t$ sont $<n$, en faisant l'hypothèse additionnelle que toutes les variables $<k$ sont liées. On a alors pour tout $t$ et $n$,
$$\verb|n_free t n = n_k_free t n 0|$$
et on peut facilement calculer \verb|n_k_free| inductivement sur $t$, en utilisant les relations suivantes :
\end{paragraph}
$$\begin{array}{rcl}
\verb€n_k_free (Var x) n k €&=&\verb€ (x<n)||(x<k)€\\
\verb€n_k_free (Fun t) n k €&=&\verb€ n_k_free t n (k+1)€\\
\verb€n_k_free (Appl t1 t2) n k €&=&\verb€ (n_k_free t1 n k) && (n_k_free t2 n k)€
\ea$$
\begin{paragraph}{}
On a ensuite montré que pour tout \verb|n|, \verb|n_free t n -> n_free t (n+1)| en démontrant tout d'abord cette propriété sur \verb|n_k_free|. La définition du prédicat ``Terme clos" peut alors simplement s'écrire \verb| closed t := n_free t 0|.
\end{paragraph}
\end{sous-partie}
\begin{sous-partie}{Substitution dans un terme}
\begin{sous-sous-partie}{Substitution unique}
\begin{paragraph}{}
On se donne un terme \verb|t|, un entier \verb|n| et un autre terme \verb|u|. On cherche à substituer \verb|u| à la variable \verb|n| dans \verb|t|, ce pour chaque occurrence libre de \verb|n|.
\end{paragraph}
\begin{paragraph}{}
Là encore, cette fonction est définie par induction sur \verb|t|. Les cas 1 et 3 sont simples à traiter :
\end{paragraph}
$$\begin{array}{rcl}
\verb€substitute_one (Var x) n u €&=&\begin{array}[t]{|cl}
\verb€ u€&\mbox{si }\verb|n=x|\\
\verb€ Var x€&\mbox{sinon}
\ea\\
\verb€substitute_one (Appl t1 t2) n u €&=&\verb|Appl (substitute_one t1 n u)| \\
&&\hspace{1.5cm}\verb|(substitute_one t2 n u)|\\
\ea$$
\begin{paragraph}{}
Le cas où \verb|t=Fun t1| est plus complexe à gérer. En effet, si 0 est une variable libre de \verb|u|, alors si on substitue directement \verb|u| à \verb|n| dans  \verb|t1|, on capture la variable libre 0 de \verb|u| par ce $\lambda$ additionnel. Pour préserver la liberté des variables libres de \verb|u|, une méthode consiste à augmenter de 1 la valeur de toutes les variables libres de \verb|u| avant la substitution. C'est la méthode que nous avons choisi d'utiliser.
\end{paragraph}
\begin{paragraph}{}
On appelle \verb|shift| cette opération. Pour la définir, nous avons utilisé une fonction \verb|add_n_to_free| qui ajoute une constante \verb|n| à toutes les variables libres d'un terme, que l'on a définit par induction sur la taille des terme, de façon similaire à la définition de \verb|n_free| (on utilise une fonction auxiliaire où on fait une hypothèse supplémentaire sur les variables liées). On a alors :
\end{paragraph}
$$\verb€substitute_one (Fun t1) n u = Fun (substitute_one t1 n (shift u))€$$
\begin{paragraph}{}
Nous avons alors pu démontrer formellement grâce à Coq que toute substitution dans un terme fermé $t$ renvoie ce même terme.
\end{paragraph}
\end{sous-sous-partie}
\begin{sous-sous-partie}{Substitution multiple}
\begin{paragraph}{}
On cherche maintenant à formaliser la substitution $t[i\leftarrow u_0\dots u_n]$, où $t, u_0, \dots, u_n$ sont des termes et $i, \dots, i+n-1$ des variables libres de $t$. Cette substitution est définie par récurrence en utilisant :
$$t[i\leftarrow u_0\dots u_n] = (t[i\leftarrow u_0])[i+1\leftarrow u_1,\dots,u_n]$$
Nous avons alors prouvé en utilisant Coq que :
$$\begin{array}{rcl}
t[] &=& t\\
t[i\leftarrow u] &=& t\quad\mbox{si $t$ satisfait $C[i](t)$}
\ea$$
Nous n'avons à l'heure où nous écrivons ce rapport pas encore démontré que
$$t[i\leftarrow u_0\dots u_n] = t[i\leftarrow u_0\dots u_n][i\leftarrow u_0] \quad\mbox{si $\forall k\le i, C[i](u_k)$}$$
mais espérons le faire avant la présentation orale du mardi 18 décembre 2012.
\end{paragraph}
\end{sous-sous-partie}
\end{sous-partie}

\end{partie}

\glsaddall
\printglossary[title={Glossaire\markright{GLOSSAIRE}}]
\addcontentsline{toc}{section}{Glossaire}


\end{document}